\section{Experimental Details}
\label{app:experimental-details}

\subsection{Models and Distillation Corpus}

We evaluate Qwen3-Embedding-4B $\rightarrow$ BERT-base (109M), BGE-M3 $\rightarrow$ MiniLMv2-H768 (66M), and Qwen3-Embedding-0.6B $\rightarrow$ MiniLMv2-H384 (22M). The unlabeled corpus contains approximately $15{,}000$ sentences, sampled evenly from the training splits of Emotion, WiC, and STS-B. Corpus texts are deduplicated before graph construction. Teacher embeddings are unit-normalized, cached once, and reused by every method variant that requires teacher representations.

\subsection{Optimization and Reproducibility}

\begin{table}[t]
\centering
\small
\setlength{\tabcolsep}{5pt}
\begin{tabular}{ll}
\toprule
Component & Configuration \\
\midrule
Optimizer & AdamW \\
Batch size and duration & $32$ anchors; $5$ epochs \\
Learning-rate schedule & $2\times10^{-5}$; $6\%$ warm-up; cosine decay \\
Minimum learning rate & $3\times10^{-6}$ \\
Random seeds & $42$, $43$, $44$ \\
Checkpoint selection & Validation average \\
\bottomrule
\end{tabular}
\caption{Training and selection protocol used for RIPPLE.}
\label{tab:optimization-config}
\end{table}

Stochastic candidate and walk streams are seeded independently by the run seed, epoch, anchor index, and stream identifier. Results for RIPPLE and TALAS are reported as mean $\pm$ standard deviation across three seeds. Comparisons of step time and peak memory use the same device, batch, mixed-precision policy, software stack, and measurement procedure for every variant; the exact accelerator and library versions should be included with the released run manifest.

\subsection{Downstream Protocols}

The nine evaluation datasets are grouped into classification (Banking77, Tweet, Emotion), pair classification (MRPC, SciTail, WiC), and semantic textual similarity (SICK, STS12, STS-B). For classification, a logistic-regression probe is fit on frozen embeddings and evaluated by macro-F1. Pair classification uses cosine similarity and is evaluated by average precision; any decision threshold used for auxiliary classification metrics is selected on validation data. STS uses the Spearman correlation between cosine similarity and human similarity scores. Emotion, WiC, and STS-B are treated as in-domain because their training data contribute unlabeled corpus sentences; the other six datasets are out-of-domain.

Checkpoints are selected using validation data, after which the corresponding test split is evaluated once. The primary aggregate is the unweighted mean of the nine task metrics. No downstream labels enter distillation or teacher-graph construction.

\subsection{Baselines}

The comparison includes the frozen teacher and undistilled student references, SimCSE-unsup, CDM, DSKD, Jasper--Stella, DistillCSE, EMO, and TALAS. Published baseline results are taken from the shared TALAS protocol where indicated in Table~\ref{tab:main_results}; TALAS and RIPPLE are summarized over three runs. Controlled RIPPLE comparisons use the same initialization, optimizer, schedule, and evaluation code and vary only the named component.

